\section{相关工作}\label{sec:related}

\textbf{可认证层与 SDP 路线。} SDPRLayers \cite{holmes2025sdpr} 将齐次 QCQP 提升为 $X\succeq0$，求解 Shor 松弛，并用一个容忍冗余约束的隐式选择定理对原始 KKT 系统求导。其已发布的位姿估计器使用 $13$ 维向量 $x=[h;\mathrm{vec}(C);t]$ 与 24 个二次约束（6 个正交性、9 个手性、9 个范数），其中最后 18 个被标为冗余。我们的工作在四个方面不同。(i) \emph{变量：}我们在 $q\in\mathbb{S}^3$（四个变量、一个约束）上优化并精确消去 $t$，因此非凸问题的维数是常数，与对应点数目无关。(ii) \emph{全局机制：}QPEP 枚举所有 KKT 根并比较目标值；不存在凸松弛，因而没有紧性假设。(iii) \emph{反向传播：}我们对约化后的平稳性系统用 $3\times3$ 黎曼 Hessian 求导；SDPRLayers 则是在行约化之后对一个带冗余行的 KKT 系统求导。(iv) \emph{歧义：}我们返回模式上的平滑后验；SDPRLayers 返回单个可认证的解，且在松弛不紧时没有定义良好的行为。我们\emph{不}主张我们的 $10\times10$ 分析适用于已发布的 24 约束实现——后者的约束张成空间是 17 维；它确实适用的完整 20D $13\times13$ 松弛，已在定理~\ref{thm:thirteen} 中证明等价于 $10\times10$ 的那一个。

\textbf{概率位姿估计。} EPro-PnP \cite{chen2022epropnp} 在 $SE(3)$ 上定义玻尔兹曼分布，并用自适应多重重要性采样以 KL 损失训练。我们的后验也是玻尔兹曼型，但支撑集是 $\mathbb{S}^3/\{\pm1\}$ 而不是 $SE(3)$，平移被精确边缘化，旋转积分是在穷举的模式列表上解析近似而非采样。这里的取舍是以偏差换方差：EPro-PnP 渐近精确但方差高，而我们的层是确定性的，但在远离分岔处带有 $O(\tau)$ 的 Laplace 偏差。公平的比较必须区分这两个误差来源，而第~\ref{sec:experiments} 节并不声称已做到这一点：只测量了离散模式采样的方差。

\textbf{QPEP 求解器。} Wu \emph{等} \cite{wu2022qpep} 把一大类位姿问题统一为单位四元数中的二次问题，消去平移，用 Gr\"obner 基求解所得多项式系统，并分析可解性、可观测性与协方差。其求解器在已发布实现中至多返回 40 个候选复根（通用的 40 根计数由第~\ref{sec:rootcount} 节的总度同伦确认）；全局解通过比较目标值获得，而不是由凸证书给出。QPEP 没有提供对数据的导数。我们的贡献正是那个缺失的可微层：闭式欧拉乘子、$3\times3$ 隐式微分、对映不变性，以及枚举根上的 Laplace 后验。我们依靠 QPEP 完成前向枚举，并把其完备性作为显式假设（A2）处理。

\textbf{局部求解器的隐式微分。} BPnP \cite{chen2020bpn} 通过对其最优性条件做隐式微分，从而穿过 PnP 求解器反向传播。该求解器是局部的，并由网络预测初始化；所得梯度是局部解的梯度，而不是全局解的梯度。在分支交叉处，局部解可能是错误的那一个，而全局 argmin 的导数作为经典函数并不存在。Theseus \cite{pineda2022theseus} 提供一个可微的非线性最小二乘库，支持李群、稀疏求解器、批处理以及隐式/直接损失最小化；它是一个通用局部优化器，不枚举全局模式，也不提供证书。我们的层更窄（四元数中的四次型代价），但它是穷举的，且反向传播为恒定规模的闭式。

\textbf{多项式优化、矩与平方和。} 我们分析的松弛是二阶 Shor/Veronese 松弛，而不是完整的矩层级。Lasserre 的层级 \cite{lasserre2001} 与 Laurent 的综述 \cite{laurent2009} 给出了一般的矩/SOS 框架；Hilbert \cite{hilbert1888} 确立了非负形式不必是平方和，Blekherman \cite{blekherman2006} 量化了非负形式比 SOS 形式多出多少。Choi--Lam 四次型 \cite{choi1977extremal,choi1980zeros} 是该间隙的标准显式见证，这也是我们用它来构造一个带认证松弛间隙的实例的原因；定理~\ref{thm:lift} 背后的 Veronese 簇的代数几何是标准的 \cite{harris1992}。我们要强调我们\emph{不}做的等同：我们的 SDP 不是教科书上带球面定位矩阵的二阶 Lasserre 矩松弛，我们也不对更高阶层级作任何主张。

\textbf{Laplace 渐近与 $SO(3)$ 的几何。} Laplace 近似及其一致（Airy）对应物是经典的渐近展开结果 \cite{bleistein1986,wong1989}；第~\ref{sec:discussion} 节的分支合并分析正是标准展开不再一致的那个退化情形。旋转积分中的测度常数来自 $SO(3)$ 的 Haar 测度与双覆盖 $\mathbb{S}^3\to SO(3)$ \cite{chirikjian2012}。定理~\ref{thm:implicit}--\ref{thm:cert} 背后的隐函数理论是 Dontchev 与 Rockafellar 的理论 \cite{dontchev2009}，而前向枚举的数值方法（阻尼牛顿、Levenberg--Marquardt、KKT 系统）是标准的 \cite{nocedal2006}。

\textbf{相对于伴随线的范围。} 伴随线 \textsc{Constrained-QPEPLayer} 通过增广拉格朗日与 Fischer--Burmeister 归约加入不等式约束（手性、非穿透、视锥）；\textsc{Robust-QPEPLayer} 加入抗外点的核以及 IRLS 权重的曲率修正。本层两者都不假设：其 $M(\theta)$ 可以是任意加权 Gram 矩阵，其理论也不依赖其他工作线中尚未发表的结果。
